\section{Introduction}
\label{sec:intro}

Ice hockey is among the least predictable of the major professional sports.
Goals are rare, around three per team per game, and a single deflection can
decide a result, so even a large difference in team quality translates into a
modest difference in win probability \citep{bennaim2006parity,lopez2018often}. For the
National Hockey League (NHL) the consequence is a narrow band of achievable
skill. On our tuning seasons a constant home-win rate scores a log loss of
0.6887 per game and a tuned Elo system \citep{elo1978rating,hvattum2010using}
scores 0.6759. Betting markets, the usual reference for the best achievable
forecast \citep{strumbelj2014determining}, are better still; our plans worked
with a rough closing-line figure of 0.66, which we could not verify because
the project uses no market data. The competition between forecasting methods
therefore takes place inside a few hundredths of a nat, and a genuine
improvement over a good rating system is a few thousandths.

Effects that small are fragile. Repeatedly trying configurations against the
same held-out seasons will eventually produce one that looks better by chance
\citep{simmons2011false,gelman2014statistical}, and information that leaks from
the future into features produces gains that vanish in use
\citep{kaufman2012leakage,kapoor2023leakage}. A model with no skill at all
would clear a 0.002-nat improvement bar on our tuning window by chance about
23\% of the time (Section~\ref{sec:protocol}). Preregistration, the practice of
committing hypotheses, analyses and decision rules before the data that test
them exist, is now standard in parts of the experimental sciences
\citep{nosek2018preregistration}, and templates for preregistering predictive
models have been proposed \citep{hofman2023preregistration}, but it remains
rare in the predictive modelling of sport.

This paper reports \neurhl{}, a system that forecasts NHL games, skaters and
seasons, built across four successive preregistrations. Each plan was committed
to the project's repository before any number that tested it was computed, and
the repository, with its full history, is public. Seasons
were assigned to development, tuning, confirmation and live roles before any
model was fitted, and the code refuses to score across those boundaries. Every
configuration searched was logged against a declared cap. Each confirmatory
test was run once, after a freeze commit, by a script that refuses a second
run. The record of what failed is kept as carefully as the record of what
worked.

Our contributions are:
\begin{enumerate}
  \item \textbf{A preregistration protocol for sports forecasting} with season
  windows, guarded data loaders, causality audits that corrupt the future and
  require the past to be bit-identical, capped search ledgers, one-shot tests,
  and a live scoring rule under which a forecast counts only if it is publicly
  committed before its game starts (Section~\ref{sec:protocol}).
  \item \textbf{A confirmed game-level improvement over Elo.} \neurhlH{} conditions
  on the dressed lineup through a neural player-level model and combines the
  result with Elo in a four-feature logistic head. On 10,184 held-out games it
  improves log loss by 0.00478 (95\% CI 0.00306 to 0.00650), in all eight
  seasons, and increases Elo's skill over the home-win base rate by about a
  quarter (Section~\ref{sec:res-h}).
  \item \textbf{A confirmed player-game model} of ice time, shots, goals and
  assists that beats each skater's own recent average on all four targets
  (Section~\ref{sec:res-pg}).
  \item \textbf{\neurhlG{}, a single-game simulation engine} that produces a full
  box score from the two lineups. Ice time is conserved exactly by
  construction, player outputs sum to team totals, and the outcome comes from a
  differentiable integration of a score-and-time hazard over the goal
  differential. It beats Elo on its gate seasons but adds no demonstrable
  win-probability information beyond \neurhlH{}, so its sealed test was not spent
  (Section~\ref{sec:res-g}).
  \item \textbf{A record of null results and failure modes}
  (Sections~\ref{sec:res-nulls} and~\ref{sec:failures}): pooled neural event
  embeddings that carry no team-strength signal; tree ensembles and wide
  networks that do not beat Elo; two leaks that passed every realism gate; a
  recording drift that changed what a shot location means; and implementation
  errors in our own gates, one of them found while writing this paper.
\end{enumerate}

Section~\ref{sec:related} places the work in context, Section~\ref{sec:data}
describes the data, Section~\ref{sec:protocol} the protocol and
Section~\ref{sec:models} the models. Section~\ref{sec:results} reports results,
Section~\ref{sec:failures} the failures and corrections,
Section~\ref{sec:live} the prospective test running during the 2026--27
season, and Section~\ref{sec:discussion} what we take from all of it. We name
seasons by the year in which they end, so 2018 is the 2017--18 season.
